\documentclass[11pt]{article}

\usepackage[utf8]{inputenc}
\usepackage[T1]{fontenc}
\usepackage[spanish]{babel}
\usepackage[a4paper,margin=2cm]{geometry}
\usepackage{xcolor}
\usepackage{listings}

\lstset{
    backgroundcolor=\color{gray!10},
    basicstyle=\ttfamily\small,
    breaklines=true,
    columns=fullflexible,
    frame=single,
    showstringspaces=false
}

\begin{document}

\begin{center}
    {\LARGE\textbf{Manual de ejecución}}\\
    \vspace{0.3cm}
    {\large Procesamiento paralelo de imágenes con Scheme y Erlang}
\end{center}

\section*{Requisitos}

El proyecto requiere las siguientes herramientas:

\begin{itemize}
    \item Erlang Base con los comandos \texttt{erl} y \texttt{erlc}
    \item Racket
    \item Full Swindle
\end{itemize}

Erlang Base y Racket se instalan con:

\begin{lstlisting}
sudo apt update
sudo apt install erlang-base racket
\end{lstlisting}

Full Swindle se instala con:

\begin{lstlisting}
raco pkg install swindle
\end{lstlisting}

\section*{Ubicación del proyecto}

La ejecución se realiza desde la carpeta principal del proyecto:

\begin{lstlisting}
cd <ruta-del-proyecto>
\end{lstlisting}

Esta carpeta contiene el archivo \texttt{image\_processor} y las carpetas
\texttt{src}, \texttt{scheme} y \texttt{samples}.

\section*{Formato de la imagen}

La imagen de entrada debe cumplir las siguientes condiciones:

\begin{itemize}
    \item Formato PPM P3
    \item Valor máximo de color igual a 255
    \item Componentes RGB entre 0 y 255
    \item Sin comentarios
\end{itemize}

Las rutas de entrada y salida deben ser diferentes.

\section*{Sintaxis de ejecución}

El archivo \texttt{image\_processor} compila el código Erlang antes de cada
ejecución.

\subsection*{Difuminado gaussiano predeterminado}

\begin{lstlisting}
./image_processor <imagen-de-entrada> <imagen-de-salida> <cantidad-de-procesos>
\end{lstlisting}

Esta forma utiliza un kernel de tamaño 3 y un sigma de 0.8493.

\begin{itemize}
    \item \verb|<imagen-de-entrada>| es la ruta de la imagen PPM P3
    \item \verb|<imagen-de-salida>| es la ruta donde se escribe el resultado
    \item \verb|<cantidad-de-procesos>| es un entero positivo
\end{itemize}

\subsection*{Difuminado gaussiano personalizado}

\begin{lstlisting}
./image_processor <imagen-de-entrada> <imagen-de-salida> \
    <cantidad-de-procesos> gaussian <tamano-del-kernel> <sigma>
\end{lstlisting}

\begin{itemize}
    \item \verb|<tamano-del-kernel>| es un número impar mayor o igual que 3
    \item \verb|<sigma>| es un número positivo
\end{itemize}

\subsection*{Escala de grises}

\begin{lstlisting}
./image_processor <imagen-de-entrada> <imagen-de-salida> \
    <cantidad-de-procesos> grayscale
\end{lstlisting}

Los nombres \texttt{gaussian} y \texttt{grayscale} se escriben en minúscula.
La cantidad de procesos no puede superar la altura de la imagen.

\section*{Cantidad de procesos}

La cantidad de procesos se cambia en el argumento
\verb|<cantidad-de-procesos>|.

Ejemplos con 1, 2, 4 y 8 procesos:

\begin{lstlisting}
./image_processor samples/sample_1.ppm sample_1_process_1.ppm 1
./image_processor samples/sample_1.ppm sample_1_process_2.ppm 2
./image_processor samples/sample_1.ppm sample_1_process_4.ppm 4
./image_processor samples/sample_1.ppm sample_1_process_8.ppm 8
\end{lstlisting}

\section*{Ejemplos}

Difuminado gaussiano predeterminado con cuatro procesos:

\begin{lstlisting}
./image_processor samples/sample_1.ppm sample_1_gaussian.ppm 4
\end{lstlisting}

Difuminado gaussiano con un kernel de tamaño 5 y un sigma de 0.75:

\begin{lstlisting}
./image_processor samples/sample_1.ppm sample_1_gaussian_5.ppm \
    4 gaussian 5 0.75
\end{lstlisting}

Escala de grises con cuatro procesos:

\begin{lstlisting}
./image_processor samples/sample_1.ppm sample_1_grayscale.ppm 4 grayscale
\end{lstlisting}

\section*{Resultado}

Por ejemplo, una ejecución correcta puede mostrar:

\begin{lstlisting}
Completed in 2.379137 seconds.
\end{lstlisting}

La imagen procesada queda en la ruta indicada por
\verb|<imagen-de-salida>|.

\end{document}
